\documentclass[../../../include/open-logic-section]{subfiles}

\begin{document}

\olfileid[id]{his}{bio}{rob}

\olsection{Julia Robinson}

\olphoto{robinson-julia}{Julia Robinson}

Julia Bowman Robinson adalah matematikawan Amerika. Ia terutama dikenal
karena karyanya tentang masalah keputusan, dan yang paling terkenal karena
sumbangannya terhadap penyelesaian masalah kesepuluh Hilbert. Robinson lahir
di St.~Louis, Missouri, pada 8 Desember 1919. Robinson mengenang bahwa ia
sudah tertarik pada bilangan sejak kecil \citep[4]{Reid1986}. Pada usia
sembilan tahun ia terserang demam scarlet dan mengalami beberapa serangan
demam rematik yang berulang. Keadaan ini memaksanya menghabiskan banyak waktu
di tempat tidur sehingga pendidikannya tertinggal. Meskipun ia dapat mengejar
ketertinggalan dengan bantuan tutor privat, dampak fisik penyakitnya
memengaruhi seluruh hidupnya.

Terlepas dari kesulitannya semasa kecil, Robinson lulus sekolah menengah
dengan beberapa penghargaan dalam matematika dan ilmu pengetahuan. Ia
memulai pendidikan universitas di San Diego State College, lalu pindah ke
University of California, Berkeley, pada tahun terakhirnya. Di sana ia
dipengaruhi oleh matematikawan Raphael Robinson. Mereka menjadi sahabat
dekat dan menikah pada tahun 1941. Sebagai istri seorang anggota fakultas,
Robinson dilarang mengajar di departemen matematika Berkeley. Meskipun tetap
mengikuti kuliah matematika sebagai pendengar, ia berharap dapat
meninggalkan universitas dan membangun keluarga. Namun, tidak lama setelah
pernikahannya, Robinson terserang pneumonia. Ia diberi tahu bahwa terdapat
penumpukan jaringan parut yang cukup besar pada jantungnya akibat demam
rematik yang dideritanya semasa kecil. Karena jaringan parut itu begitu
parah, dokter memperkirakan bahwa ia tidak akan hidup melewati usia empat
puluh tahun dan menyarankannya untuk tidak mempunyai anak
\citep[13]{Reid1986}.

Robinson mengalami depresi untuk waktu yang lama, tetapi akhirnya memutuskan
untuk melanjutkan studi matematika. Ia kembali ke Berkeley dan menyelesaikan
PhD pada tahun 1948 di bawah bimbingan Alfred Tarski. Tarski telah
menunjukkan bahwa teori orde pertama bilangan real dapat diputuskan, dan
dari karya G\"odel diperoleh bahwa teori orde pertama bilangan asli tidak
dapat diputuskan. Masalah terbuka yang besar ketika itu adalah apakah teori
orde pertama bilangan rasional dapat diputuskan. Dalam tesisnya
\citeyearpar{Robinson1949}, Robinson membuktikan bahwa teori tersebut tidak
dapat diputuskan.

Karena tertarik pada masalah keputusan, Robinson selanjutnya berusaha
menemukan penyelesaian bagi masalah kesepuluh Hilbert. Masalah ini merupakan
salah satu dari daftar terkenal 23 masalah matematika yang diajukan David
Hilbert pada tahun 1900. Masalah kesepuluh menanyakan apakah ada algoritma
yang dapat menentukan, dalam waktu berhingga, apakah suatu persamaan
polinomial berkoefisien bulat, seperti $3x^2 - 2y +3 = 0$, mempunyai solusi
dalam bilangan bulat. Pertanyaan seperti itu dikenal sebagai \emph{masalah
Diophantine}. Setelah meraih beberapa keberhasilan awal, Robinson bergabung
dengan Martin Davis dan Hilary Putnam, yang juga mengerjakan masalah
tersebut. Mereka berhasil menunjukkan bahwa masalah Diophantine eksponensial
(yang peubahnya juga boleh muncul sebagai eksponen) tidak dapat diputuskan,
dan menunjukkan bahwa suatu konjektur tertentu (yang kemudian disebut
``J.R.'') mengakibatkan masalah kesepuluh Hilbert tidak dapat diputuskan
\citep{DavisPutnamRobinson1961}. Robinson terus mengerjakan masalah tersebut
sepanjang dasawarsa 1960-an. Pada tahun 1970, matematikawan muda Rusia Yuri
Matijasevich akhirnya membuktikan hipotesis J.R. Hasil gabungan itu sekarang
disebut teorema Matijasevich--Robinson--Davis--Putnam, atau secara singkat
teorema MRDP itu. Matijasevich dan Robinson menjadi sahabat dan berkolaborasi dalam
beberapa makalah. Dalam sebuah surat kepada Matijasevich, Robinson pernah
menulis bahwa ``sebenarnya saya sangat senang bahwa dengan bekerja bersama
(terpisah ribuan mil), kita jelas mencapai kemajuan yang lebih besar daripada
yang dapat dicapai salah seorang dari kita sendirian''
\citep[45]{Matijasevich1992}.

Robinson adalah presiden perempuan pertama American Mathematical Society dan
perempuan pertama yang terpilih menjadi anggota National Academy of
Sciences. Ia meninggal pada 30 Juli 1985 dalam usia 65 tahun setelah
didiagnosis menderita leukemia.

\begin{reading}
Makalah-makalah matematika Robinson tersedia dalam \textit{Collected Works}
\citep{Robinson1996}, yang juga memuat cetak ulang memoir biografisnya dari
National Academy of Sciences \citep{Feferman1994}. Kakak perempuan Robinson,
Constance Reid, menerbitkan ``Autobiography of Julia,'' berdasarkan
wawancara \citep{Reid1986}, serta sebuah memoir lengkap
\citep{Reid1996}. Sebuah dokumenter pendek mengenai Robinson dan masalah
kesepuluh Hilbert disutradarai oleh George Csicsery \citep{Csicsery2016}.
Untuk memoir singkat mengenai kolaborasi Yuri Matijasevich dengan Robinson
dan pengaruh Robinson terhadap karyanya, lihat \citep{Matijasevich1992}.
\end{reading}

\end{document}
